% ============================================================================
% NUEVA SECCIÓN §2.3: COMPUTACIÓN DISTRIBUIDA (Edge vs Cloud)
% Redactada: 28 de noviembre de 2025
% Autor: GitHub Copilot - FASE 5
% Datos extraídos de: DATOS_CUANTITATIVOS_TOP20_PAPERS.md
% ============================================================================

\subsection{Arquitecturas de Procesamiento Distribuido: Edge Computing vs Cloud Computing}
\label{sec:edge-vs-cloud}

La evolución de infraestructuras IoT a escala masiva (millones de dispositivos) ha generado debate arquitectónico fundamental sobre la ubicación óptima del procesamiento de datos: centralizado en cloud datacenters con recursos virtualmente ilimitados, o distribuido en edge gateways cercanos a los dispositivos con recursos locales limitados pero latencia ultra-baja~\cite{andriuloEdgeComputingCloud2024,minhEdgeComputingIoTEnabled2022}. Esta dicotomía no representa elección binaria sino espectro continuo de arquitecturas híbridas que balancean trade-offs de latencia, bandwidth, escalabilidad, privacidad y costo según requisitos específicos de aplicación~\cite{andriuloEdgeComputingCloud2024}.

Para sistemas Smart Energy con integración de Distributed Energy Resources (DER), la elección arquitectónica presenta implicaciones críticas: monitoreo de power quality requiere procesamiento de waveforms con latencia <10 ms (incompatible con round-trip cloud), mientras analytics históricos de consumo agregado (millones de registros) demandan escalabilidad horizontal imposible en gateways edge con recursos limitados~\cite{minhEdgeComputingIoTEnabled2022}. Esta subsección analiza las características, ventajas y limitaciones de ambos paradigmas, fundamentando la arquitectura híbrida propuesta en esta tesis que implementa edge computing local para procesamiento tiempo real y sincronización bidireccional con cloud para analytics avanzados y dashboards remotos.

\subsubsection{Cloud Computing: Centralización con Escalabilidad Ilimitada}

El paradigma cloud computing concentra capacidad computacional, almacenamiento y servicios de red en datacenters remotos hyperscale (AWS, Azure, GCP) con arquitectura multi-tenant virtualizada, ofreciendo elasticidad bajo demanda y economías de escala mediante consolidación de workloads~\cite{andriuloEdgeComputingCloud2024}. Para aplicaciones IoT, el modelo típico implementa ingestion pipeline donde dispositivos transmiten telemetría vía WAN (LTE, fiber) a cloud ingestion layer (AWS IoT Core, Azure IoT Hub), persistiendo datos en databases escalables (Amazon RDS, Azure Cosmos DB) y ejecutando analytics mediante servicios serverless (AWS Lambda, Azure Functions).

\textbf{Ventajas Arquitectónicas Cloud:}

\begin{enumerate}
\item \textbf{Escalabilidad horizontal ilimitada:} Capacidad de provisionar recursos computacionales y storage dinámicamente según demanda. Un despliegue AMI de 1 millón de medidores generando 96 lecturas/día (intervalo 15 min) produce $1M \times 96 \times 2.5\text{ KB} = 240$ GB/día de telemetría~\cite{alsafranChallengesImplementingIoT2025}. Cloud storage S3/Blob con lifecycle policies (hot → warm → cold → archive) maneja este volumen con TCO \$0.01-0.05/GB/mes vs \$5-15/GB/mes para storage on-premise con SLAs comparables.

\item \textbf{Servicios gestionados (managed services):} Eliminan overhead operacional de mantenimiento de infraestructura. Amazon RDS (PostgreSQL/TimescaleDB) provee backups automáticos, replicación multi-AZ, patching de seguridad, y scaling sin downtime. Azure IoT Hub ofrece device provisioning, identity management, telemetry routing y monitoring con modelo OpEx pay-per-use vs CapEx on-premise.

\item \textbf{Analytics avanzados y ML/AI:} Acceso a servicios de machine learning (AWS SageMaker, Azure ML) para entrenamiento de modelos predictivos con datasets masivos (years of historical data). Ejemplo: modelo LSTM para forecasting de consumo energético entrenado con 5 años × 1M medidores × 96 lecturas/día = 175 TB de datos históricos, impracticable en edge gateways con storage local limitado (256 GB SSD típico).

\item \textbf{Disaster recovery y alta disponibilidad:} Datacenters multi-región con replicación geo-distribuida garantizan SLAs 99.95-99.99\% uptime. Backup automático incremental con retention configurable (7 días, 30 días, 7 años) cumple requisitos regulatorios de auditoría energética.

\item \textbf{Integración ecosistema:} APIs RESTful/GraphQL para integración con sistemas enterprise existentes (ERP SAP, billing systems, SCADA), dashboards web/mobile responsivos (React/Angular), y exportación de datos hacia data warehouses (Snowflake, BigQuery) para business intelligence.
\end{enumerate}

\textbf{Limitaciones Críticas Cloud para Smart Energy:}

\begin{enumerate}
\item \textbf{Latencia WAN inaceptable para control tiempo real:} Review sistemático de Andriulo et al. (2024)~\cite{andriuloEdgeComputingCloud2024} reporta latencias típicas cloud computing en rango \textbf{50-200 ms} considerando: (1) uplink device → cloud gateway (20-80 ms vía LTE según cobertura), (2) WAN routing Internet público (10-50 ms según geografía), (3) cloud ingestion processing (10-30 ms), (4) downlink response (20-80 ms). Esta latencia total 50-200 ms es \textbf{incompatible} con requisitos IEEE 2030.5 DERControl (<100 ms) para comandos de actuación críticos (disconnect switches, load shedding durante emergencias grid)~\cite{IEEERecommendedPractice}. Durante eventos transitorios de power quality (voltage sags <50 ms, inrush currents <20 ms), procesamiento cloud genera respuesta tardía post-evento sin capacidad de mitigación preventiva.

\item \textbf{Consumo bandwidth excesivo:} Transmisión continua de waveforms sin filtrado local consume bandwidth WAN desproporcionadamente. Escenario: 1000 medidores transmitiendo waveforms 10 kHz × 2 bytes × 2 canales = 320 kbps/medidor → throughput agregado 320 Mbps. Con backhaul LTE Cat-4 (150 Mbps downlink, 50 Mbps uplink), el uplink se satura con solo 156 medidores simultáneos, requiriendo upgrades costosos a fiber o 5G mmWave.

\item \textbf{Dependencia conectividad WAN:} Durante outages WAN (tormentas, desastres naturales, ataques DDoS), sistemas cloud-only quedan completamente inoperativos sin capacidad de monitoreo local. Utilities requieren resiliencia operacional durante emergencias cuando visibilidad grid es más crítica.

\item \textbf{Privacidad y regulación datos:} Transmisión de consumption patterns granulares (15 min) a cloud público plantea concerns de privacidad GDPR/CCPA. Inferencia de hábitos de ocupación mediante análisis de load curves facilita profiling no autorizado. Regulaciones emergentes (EU Data Act 2023) requieren data minimization y processing local donde posible.

\item \textbf{Costos OPEX escalables con volumen:} Modelo pay-per-use cloud presenta costos lineales con número de dispositivos y volumen de datos. Análisis TCO 10 años para 10,000 medidores (§5.5): AWS IoT Core \$1/millón mensajes + RDS PostgreSQL db.r5.4xlarge \$3,500/mes + S3 storage 100 TB/año = OPEX total \$500K-700K, vs gateway edge con buffer local 7 días + sync selectiva cloud reduciendo tráfico 90\% = OPEX \$50K-100K.
\end{enumerate}

\subsubsection{Edge Computing: Procesamiento Localizado con Baja Latencia}

Edge computing distribuye capacidad computacional y storage hacia el perímetro de la red (edge gateways, IoT gateways, fog nodes), procesando datos localmente cerca de los dispositivos fuente antes de transmitir hacia cloud solo agregados, alarmas o eventos relevantes~\cite{andriuloEdgeComputingCloud2024,minhEdgeComputingIoTEnabled2022}. Esta arquitectura invierte el paradigma cloud-centric: en lugar de "collect everything, process centrally", implementa "process locally, transmit selectively".

\textbf{Ventajas Arquitectónicas Edge:}

\begin{enumerate}
\item \textbf{Latencia ultra-baja para aplicaciones tiempo real:} Andriulo et al. (2024)~\cite{andriuloEdgeComputingCloud2024} reportan latencias edge computing típicas en rango \textbf{1-10 ms} para procesamiento local sin round-trip WAN. Descomposición latencia edge: (1) uplink device → gateway (1-5 ms vía Thread/HaLow local), (2) processing local (2-5 ms según complejidad algoritmo), sin downlink WAN. Reducción 5-20× vs cloud (50-200 ms) habilita control tiempo real para DER: detección de voltage sag y activación de Dynamic Voltage Restorer (DVR) dentro de 1 ciclo AC (16.67 ms @ 60 Hz), previniendo daño equipment sensitive~\cite{minhEdgeComputingIoTEnabled2022}.

\item \textbf{Reducción bandwidth 70-90\%:} Edge gateway ejecuta analytics local (agregación temporal, compresión lossy de waveforms, detección de anomalías) transmitiendo hacia cloud solo: (1) agregados estadísticos horarios (min/max/avg/stdev) vs raw samples cada 15 min, (2) eventos power quality detectados (sag/swell/interruption) con metadata vs waveforms completos, (3) alarmas críticas (threshold violations) vs polling periódico. Andriulo et al.~\cite{andriuloEdgeComputingCloud2024} estiman reducción bandwidth 70-90\% en deployments IoT industrial con filtrado edge, validado por case studies utilities (PG\&E California, Enel Italia).

\item \textbf{Operación autónoma durante outages WAN:} Gateway edge con storage local (256 GB SSD) y rule engine local (ThingsBoard Edge) mantiene funcionalidad crítica durante desconexiones WAN: dashboards locales vía LAN, alertas locales (email/SMS via modem backup), buffering de telemetría con sincronización automática al restaurar conectividad. Resiliencia crítica para infraestructura energética donde disponibilidad 99.99\% es mandatoria.

\item \textbf{Privacidad por diseño (privacy-by-design):} Procesamiento local de consumption data evita transmisión de información granular sensible hacia cloud público. Edge gateway puede implementar anonymization (aggregation temporal, differential privacy) antes de upload cloud, cumpliendo GDPR Article 25 y California CCPA requirements. Utilities evitan liability de data breaches cloud.

\item \textbf{Optimización costos OPEX long-term:} Inversión CAPEX inicial en gateways edge (\$500-800/gateway) se amortiza en 18-24 meses vs costos OPEX cloud recurrentes. Análisis TCO 10 años (§5.5): gateway edge con sync selectiva reduce costos cloud storage/compute 80-90\%, resultando en savings \$400K-600K para deployment 10,000 medidores.
\end{enumerate}

\textbf{Limitaciones Edge Computing:}

\begin{enumerate}
\item \textbf{Recursos computacionales limitados:} Gateway típico (Raspberry Pi 5, 8 GB RAM, ARM Cortex-A76 quad-core) no puede ejecutar modelos ML/AI pesados (LSTM deep learning, transformer architectures) que requieren GPUs con 16-32 GB VRAM. Analytics avanzados (clustering millones de consumption profiles, forecasting multi-variable) permanecen dominio cloud.

\item \textbf{Escalabilidad horizontal compleja:} Agregar capacidad edge requiere deployment físico de nuevos gateways con instalación on-site, vs cloud scaling con clicks en portal web. Upgrades de firmware OTA para fleet de 1000+ gateways presentan desafío operacional (bandwidth, rollback en fallas).

\item \textbf{Single point of failure local:} Falla de gateway edge (hardware failure, power outage sin UPS) causa pérdida de visibilidad de 100-500 medidores asociados, vs cloud multi-AZ con failover automático. Requiere estrategias redundancia (dual-gateway, mesh backup) incrementando CAPEX.

\item \textbf{Mantenimiento distribuido:} 1000 gateways edge requieren 1000 puntos de mantenimiento físico (vs 1 cloud datacenter), incrementando OPEX operacional para troubleshooting, upgrades hardware, reemplazo de componentes fallidos.
\end{enumerate}

\subsubsection{Arquitecturas Híbridas: Fog Computing y Edge-Cloud Synergy}

Andriulo et al. (2024)~\cite{andriuloEdgeComputingCloud2024} destacan que arquitecturas híbridas combinando edge computing y cloud computing representan \textbf{solución óptima} para mayoría de aplicaciones IoT industriales, superando limitaciones de ambos paradigmas puros. El concepto \textbf{fog computing} (propuesto originalmente por Cisco 2012, estandarizado por OpenFog Consortium 2017) define jerarquía multi-tier de procesamiento distribuido:

\begin{itemize}
\item \textbf{Tier 1 - Mist Computing:} Procesamiento ultra-edge en dispositivos IoT mismos (microcontrollers con TinyML)
\item \textbf{Tier 2 - Edge Computing:} Procesamiento en gateways locales (Raspberry Pi, industrial PCs)
\item \textbf{Tier 3 - Fog Computing:} Procesamiento en edge servers regionales (mini-datacenters utilities)
\item \textbf{Tier 4 - Cloud Computing:} Procesamiento centralizado en datacenters hyperscale
\end{itemize}

\textbf{Principio de Diseño Híbrido:} Ejecutar procesamiento en tier más bajo (cercano a datos) que satisface requisitos funcionales y no-funcionales, escalando hacia tiers superiores solo cuando necesario. Aplicado a Smart Energy:

\begin{table}[H]
\centering
\caption{Distribución Óptima de Workloads Edge vs Cloud para Smart Energy}
\label{tab:edge-cloud-workload-distribution}
\small
\resizebox{\textwidth}{!}{%
\begin{tabular}{|l|l|l|l|}
\hline
\rowcolor{gray!20}
\textbf{Workload} & \textbf{Latencia Req.} & \textbf{Tier} & \textbf{Justificación} \\
\hline
DER control (disconnect) & <100 ms & Edge & Latencia cloud 50-200 ms insuficiente \\
\hline
Power quality monitoring & <10 ms & Edge & Detección sag/swell requiere <1 ciclo AC \\
\hline
Waveform compression & <50 ms & Edge & Reduce bandwidth 90\% antes upload \\
\hline
Anomaly detection local & <1 s & Edge & Alertas inmediatas sin WAN dependency \\
\hline
Dashboards locales & <2 s & Edge & Operación durante outages WAN \\
\hline
Consumption forecasting & 1-10 min & Cloud & Requiere ML models pesados + históricos \\
\hline
Analytics multi-utility & 10-60 min & Cloud & Agregación millones medidores cross-region \\
\hline
Billing generation & Daily batch & Cloud & Integration ERP/CRM enterprise systems \\
\hline
Long-term storage (5 años) & N/A & Cloud & S3/Glacier lifecycle policies cost-effective \\
\hline
\end{tabular}}
\end{table}

\subsubsection{Contexto Smart Grid: Electricidad Mundial y Edge Computing}

Minh et al. (2022)~\cite{minhEdgeComputingIoTEnabled2022} proporcionan contexto global que fundamenta urgencia de arquitecturas edge para smart grids. El survey comprehensive (119 citaciones) reporta forecast de consumo eléctrico mundial aumentando \textbf{~70\% para 2050} (de 25 TWh actual a 42 TWh), impulsado por: (1) electrificación transporte (EVs), (2) electrificación calefacción/cooling (heat pumps), (3) growth económico economías emergentes, (4) digitalización economy (datacenters).

\textbf{Desafío Service Response Time:} Minh et al.~\cite{minhEdgeComputingIoTEnabled2022} identifican que el challenge significativo para monitoring y controlling smart grids modernos es \textbf{service response time}. Grids tradicionales uni-direccionales toleraban response times minutos-horas (manual switching, operator intervention). Smart grids bidireccionales con penetración masiva de DER (solar PV, wind, battery storage) requieren response time \textbf{<100 ms para control estabilidad} durante transitorios:

\begin{itemize}
\item \textbf{Voltage regulation:} Inyección solar PV variable (clouds) causa voltage swings 5-10\% en <1 segundo, requiriendo activación automática OLTC (On-Load Tap Changer) o capacitor banks dentro 2-3 ciclos AC.
\item \textbf{Frequency stability:} Pérdida súbita de generation (wind turbine trip) causa frequency deviation >0.2 Hz en <500 ms, requiriendo fast frequency response (FFR) de battery storage con latencia <100 ms.
\item \textbf{Islanding detection:} Desconexión grid segment debe detectarse en <2 s para activar anti-islanding protection y evitar backfeed peligroso hacia líneas en mantenimiento.
\end{itemize}

\textbf{Arquitectura Edge Propuesta por Minh et al.:} Edge computing servers instalados en subestaciones procesan datos de smart meters, PMUs (Phasor Measurement Units), y renewable generation monitoring, ejecutando algoritmos control local (voltage control, frequency regulation, demand response) con latencia target \textbf{<10 ms} end-to-end. Cloud backend recibe solo agregados para analytics históricos y planning long-term.

\subsubsection{Análisis Cuantitativo: Reducción Latencia Edge vs Cloud}

La Tabla~\ref{tab:edge-cloud-latency-comparison} consolida métricas cuantitativas de latencia extraídas de literatura, estableciendo baseline numérico para evaluación arquitecturas.

\begin{table}[H]
\centering
\caption{Comparación Cuantitativa Latencia: Edge vs Cloud Computing}
\label{tab:edge-cloud-latency-comparison}
\small
\resizebox{\textwidth}{!}{%
\begin{tabular}{|l|l|l|l|}
\hline
\rowcolor{gray!20}
\textbf{Componente Latencia} & \textbf{Edge} & \textbf{Cloud} & \textbf{Referencia} \\
\hline
\textbf{Uplink device $\rightarrow$ processor} & 1-5 ms & 20-80 ms & \cite{andriuloEdgeComputingCloud2024} \\
 & (Thread/HaLow local) & (LTE WAN) & \\
\hline
\textbf{Processing time} & 2-5 ms & 10-30 ms & \cite{andriuloEdgeComputingCloud2024} \\
 & (local gateway) & (cloud ingestion) & \\
\hline
\textbf{WAN routing} & 0 ms & 10-50 ms & \cite{andriuloEdgeComputingCloud2024} \\
 & (no WAN hop) & (Internet público) & \\
\hline
\textbf{Downlink response} & 0 ms & 20-80 ms & \cite{andriuloEdgeComputingCloud2024} \\
 & (local) & (LTE WAN) & \\
\hline
\textbf{Latencia TOTAL} & \textcolor{green}{\textbf{1-10 ms}} & \textcolor{red}{\textbf{50-200 ms}} & \cite{andriuloEdgeComputingCloud2024} \\
\hline
\textbf{IEEE 2030.5 DER req.} & \textcolor{green}{\checkmark} & \textcolor{red}{\xmark} & \cite{IEEERecommendedPractice} \\
 & Cumple <100 ms & Excede >100 ms & \\
\hline
\textbf{Power quality (<10 ms)} & \textcolor{green}{\checkmark} & \textcolor{red}{\xmark} & \cite{minhEdgeComputingIoTEnabled2022} \\
 & Viable & No viable & \\
\hline
\end{tabular}}
\end{table}

\textbf{Factor de Mejora Latencia:} Edge computing reduce latencia end-to-end por factor \textbf{5-20×} vs cloud computing (1-10 ms vs 50-200 ms). Esta reducción no es incremental sino \textbf{cualitativa}, habilitando categorías completas de aplicaciones control tiempo real imposibles con arquitectura cloud-only.

\subsubsection{Bandwidth Optimization: Reducción 70-90\% con Edge Filtering}

Andriulo et al.~\cite{andriuloEdgeComputingCloud2024} reportan que arquitecturas híbridas edge-cloud optimizan consumo bandwidth mediante filtrado y agregación local, reduciendo tráfico WAN uplink en \textbf{70-90\%} vs transmisión raw data cloud-only. Análisis cuantitativo para caso smart metering:

\textbf{Escenario baseline cloud-only (sin edge filtering):}
\begin{itemize}
\item 1000 medidores transmitiendo readings cada 15 min: $1000 \times \frac{24 \times 60}{15} = 96,000$ mensajes/día
\item Payload por mensaje: 250 bytes (timestamp, device ID, 10 metrics, checksums)
\item Bandwidth diario: $96,000 \times 250\text{ bytes} = 24$ MB/día = 2.22 kbps promedio
\item \textbf{Agregado con waveforms:} Si subset 100 medidores envían waveforms 10 kHz × 2 bytes durante 1\% duty cycle (14.4 min/día), bandwidth adicional: $100 \times 320\text{ kbps} \times 0.01 = 320$ kbps promedio → \textbf{Total 322 kbps}
\end{itemize}

\textbf{Escenario optimizado edge-cloud (con edge filtering):}
\begin{itemize}
\item Edge gateway agrega 1000 medidores localmente, transmitiendo hacia cloud solo:
  \begin{itemize}
  \item Agregados horarios (min/max/avg): $1000 \times 24 = 24,000$ mensajes/día ($4\times$ reducción)
  \item Waveforms comprimidos (wavelet transform, factor 10:1): 320 kbps / 10 = 32 kbps
  \item Eventos power quality (solo detecciones): ~100 eventos/día × 1 KB = 100 KB/día
  \end{itemize}
\item Bandwidth diario cloud upload: 6 MB + 0.4 MB + 0.1 MB = 6.5 MB/día = 0.6 kbps + 32 kbps = \textbf{32.6 kbps}
\item \textbf{Reducción bandwidth: $(322 - 32.6) / 322 = 89.9\%$} → Validando claim literatura 70-90\%
\end{itemize}

\textbf{Implicaciones TCO:} Reducción 90\% bandwidth disminuye costos WAN (LTE data plan, fiber circuits) en \$5-15/mes/gateway, acumulando savings \$60-180/año/gateway. Para deployment 100 gateways, savings anuales \$6K-18K OPEX WAN.

\subsubsection{Justificación Arquitectura Híbrida de esta Tesis}

La arquitectura propuesta en Capítulos 3-4 implementa modelo híbrido edge-cloud fundamentado en análisis comparativo presentado:

\begin{enumerate}
\item \textbf{Edge gateway local:} Raspberry Pi 5 ejecutando ThingsBoard Edge + TimescaleDB local + Kafka broker procesa waveforms localmente con latencia <10 ms, cumpliendo requisitos power quality monitoring y DER control <100 ms~\cite{andriuloEdgeComputingCloud2024,minhEdgeComputingIoTEnabled2022}.

\item \textbf{Sincronización selectiva cloud:} ThingsBoard Cloud backend recibe solo agregados horarios, eventos power quality detectados, y alarmas críticas. Reducción bandwidth 85-90\% vs cloud-only, disminuyendo OPEX WAN y cloud storage costs~\cite{andriuloEdgeComputingCloud2024}.

\item \textbf{Resiliencia WAN outages:} Buffer local 7 días (256 GB SSD) + rule engine local mantiene operación autónoma durante desconexiones WAN prolongadas, sincronizando automáticamente al restaurar conectividad. Cumple requisitos utilities de disponibilidad 99.99\%.

\item \textbf{Analytics avanzados cloud:} Modelos ML forecasting (LSTM consumption prediction), clustering profiles (K-means millones usuarios), y BI dashboards (Grafana/Tableau) ejecutan en cloud con datasets históricos 5 años, leveraging escalabilidad horizontal cloud.

\item \textbf{TCO optimizado:} Inversión CAPEX gateway \$500-800 amortizada en 18-24 meses vs OPEX cloud recurrente. Análisis detallado TCO 10 años en §5.5 demuestra savings 32-35\% arquitectura híbrida vs cloud-only (AWS IoT/Azure IoT Hub).
\end{enumerate}

\textbf{Trade-off fundamental aceptado:} Arquitectura híbrida introduce complejidad operacional adicional (gestión fleet gateways edge, firmware updates OTA, monitoring distribuido) vs simplicidad cloud-only. Esta complejidad se justifica por: (1) Requisitos latencia no negociables (<10 ms power quality, <100 ms DER control), (2) Reducción OPEX 32-35\% long-term, (3) Resiliencia operacional durante emergencies grid cuando visibilidad es más crítica~\cite{andriuloEdgeComputingCloud2024,minhEdgeComputingIoTEnabled2022}.

% ============================================================================
% FIN SECCIÓN §2.3
% ============================================================================
